% main_bsiMC_ufu.tex v1.1, Thiago Pirola Ribeiro
% adaptado do modelo PPGCO - main_ppgco_ufu.tex v1.0, Lásaro Camargos e Denise Guliato
% Criacao: 2019-03-08 - v1.0
% ------------------------------------------------------------------------
% Atualizações.
% 2022-10-18 - v1.1 - Thiago Pirola Ribeiro
% 2025-03-12 - V1.2 - Diego Nunes Molinos
%          - Atualização do termo grau para Título de Bacharel;
%          - Inclusão da seção de Abstract; 
% ------------------------------------------------------------------------
\documentclass{bcc}
%Não altere o comando seguinte. O título de seu trabalho será especificado mais adiante.
\title{Análise e Predição de Trajetórias de Espermatozoides em Vídeos Microscópicos}

% ---
% Pacotes fundamentais 
% ---
%\usepackage{cmap}				% Mapear caracteres especiais no PDF (já carregado pela classe)
\usepackage{lmodern}				% Usa a fonte Latin Modern			
\usepackage{makeidx}            	% Cria o indice
\usepackage{hyperref}  			% Controla a formação do índice
\usepackage{lastpage}			% Usado pela Ficha catalográfica
\usepackage{indentfirst}			% Indenta o primeiro parágrafo de cada seção.
\usepackage{nomencl} 			% Lista de simbolos
\usepackage{graphicx}			% Inclusão de gráficos
% ---

% ---
% Pacotes adicionais, usados apenas no âmbito do Modelo eesc
% ---
\usepackage[printonlyused]{acronym}
\usepackage[table]{xcolor}
\usepackage{amsmath}
\usepackage{booktabs}
\usepackage{tabularx}
\usepackage{array}
\usepackage{float}
\newcolumntype{Y}{>{\raggedright\arraybackslash}X}
% ---

% ---
% Insira aqui os pacotes utilizados por você
% ---


% ---
% Informações de dados para CAPA e FOLHA DE ROSTO
% ---
%
% Título:
%	1. Título em português
%	2. Título em inglês
\titulo
{Análise e Predição de Trajetórias de Espermatozoides em Vídeos Microscópicos}
{Analysis and Prediction of Spermatozoa Trajectories in Microscopy Videos}

% Colocar aqui ate 5 palavras-chave do trabalho
\keyword{visão computacional}
\keyword{espermatozoides}
\keyword{rastreamento multiobjeto}
\keyword{fluxo óptico}
\keyword{predição de trajetórias}

%
% Autor:
%	1. Nome completo do autor
%	2. Formato de nome para bibliografia
\autor{Leandro Marques Gontijo Jersé}{Jersé, Leandro Marques Gontijo}
% Cidade
\local{Uberlândia - MG}
%
% Ano de defesa
\date{2026}
% Área de concentração da pesquisa
\areaconcentracao{Ciência da Computação}
% Nome do orientador
\orientador{Prof. Dr. Bruno Augusto Nassif Travencolo}
% Nome do coorientador, se houver
%\coorientador{Nome completo do coorientador}
% -


% informações do PDF
\hypersetup{
	pdftitle={\imprimirtitulo},
	pdfauthor={\imprimirautor},
	pdfsubject={\imprimirpreambulo},
%	pdfkeywords=\@resumokw,
	pdfproducer={LaTeX with abnTeX2}, 	% producer of the document
	colorlinks=true,       		% false: boxed links; true: colored links
	linkcolor=black,          	% color of internal links
	citecolor=black,        		% color of links to bibliography
	filecolor=magenta,      		% color of file links
	urlcolor=black,
	bookmarksdepth=4
}

% ---
% compila o indice
% ---
\makeindex
% ---

% ---
% Compila a lista de abreviaturas e siglas
% ---
\makenomenclature
% ---

% ----
% Início do documento
% ----

\begin{document}

	\nomenclature{$\mathbf{p}_t$}{posição bidimensional de uma célula no quadro $t$}
	\nomenclature{$\Delta t$}{intervalo de tempo entre dois quadros consecutivos}
	\nomenclature{$T$}{quantidade de posições de uma trajetória}
	\nomenclature{$u,v$}{componentes horizontal e vertical do fluxo óptico}
	\nomenclature{$p$}{quantidade de quadros da janela observada}
	\nomenclature{$q$}{quantidade de quadros do horizonte de predição}
	
	% ----------------------------------------------------------
	% ELEMENTOS PRÉ-TEXTUAIS
	% ----------------------------------------------------------
	\pretextual
	
	% ---
	% Insere Capa, Folha de rosto, Ficha catalográfica (se inserida)
	% e folha de aprovação (se inserida).
	% ---
	\maketitle

	% ---
	% Dedicatória - obrigatório somente para TCC 2
	% ---
	%\imprimirdedicatoria{Texto da dedicatória}
	% ---
	
	% ---
	% Agradecimentos - obrigatório somente para TCC 2
	% ---
	%\imprimiragradecimentos{Texto dos agradecimentos}
	% ---
	
	% ---
	% Epígrafe - obrigatório somente para TCC 2
	% ---
	%\imprimirepigrafe{Texto da epígrafe}
	% ---
	
	% ---
	% RESUMO - obrigatório somente para TCC 2
	% ---
	
	% Resumo em português
	\begin{resumo}{}
		Este trabalho investiga a análise e a predição de trajetórias de espermatozoides em vídeos microscópicos por meio de técnicas de Visão Computacional. O problema envolve objetos pequenos, visualmente semelhantes e de baixo contraste, cujos movimentos rápidos, irregulares e sujeitos a sobreposições dificultam a detecção e a preservação de identidades ao longo do tempo. O objetivo é implementar e avaliar uma pipeline composta por detecção, rastreamento multiobjeto, estimação de movimento aparente por fluxo óptico e predição de posições futuras. A pesquisa será aplicada, quantitativa e experimental. Serão utilizados exclusivamente vídeos microscópicos do conjunto público VISEM, com as anotações do VISEM-Tracking disponíveis para trechos de 20 desses vídeos. As divisões experimentais serão independentes por vídeo de origem, mantendo todos os trechos de um mesmo participante na mesma divisão. Métodos clássicos, como limiarização, Filtro de Kalman, SORT e Farnebäck, serão comparados a abordagens aprendidas, como YOLO, ByteTrack, RAFT e redes recorrentes. A hipótese principal é que características locais de fluxo óptico, combinadas ao histórico de posições, podem reduzir os erros ADE e FDE em relação a preditores equivalentes que utilizem somente posição e velocidade. A avaliação também utilizará métricas de detecção e rastreamento, com atenção à reprodutibilidade e à distinção entre movimento aparente na imagem e velocidade física do fluido.
	\end{resumo}	

    % Abstract
	\begin{abstract}{Computer vision. Spermatozoa. Multi-object tracking. Optical flow. Trajectory prediction}
		This work investigates the analysis and prediction of spermatozoa trajectories in microscopy videos using Computer Vision techniques. The problem involves small, visually similar, low-contrast objects whose fast and irregular motion, overlaps, and temporary occlusions make detection and identity preservation difficult. The objective is to implement and evaluate a pipeline composed of detection, multi-object tracking, apparent-motion estimation through optical flow, and future-position prediction. The study will follow an applied, quantitative, and experimental design. Microscopy videos will come exclusively from the public VISEM dataset, with VISEM-Tracking annotations available for segments of 20 of these videos. Experimental splits will be independent by source video, keeping all segments from the same participant in the same split. Classical methods, including thresholding, Kalman filtering, SORT, and Farnebäck optical flow, will be compared with learned approaches such as YOLO, ByteTrack, RAFT, and recurrent neural networks. The main hypothesis is that local optical-flow features combined with position histories can reduce ADE and FDE relative to equivalent predictors that use only position and velocity. Detection and tracking metrics will also be reported, with emphasis on reproducibility and on the distinction between image-based apparent motion and the physical velocity of the fluid.
	\end{abstract}
		
	% ---
	% inserir lista de ilustrações
	% ---
	%\listailustracoes
	% ---
	
	% ---
	% inserir lista de tabelas
	% ---
	\listatabelas
	% ---
	
	% ---
	% inserir lista de abreviaturas e siglas
	% ---
	\listasiglas{abrev/Abreviaturas}
	% ---

	% ---
	% inserir lista de símbolos
	% ---
	\listasimbolos
	% ---
	
	% ---
	% inserir o sumario
	% ---
	\sumario
	% ---
	
	% ----------------------------------------------------------
	% ELEMENTOS TEXTUAIS
	% ----------------------------------------------------------
	\mainmatter
	
	
	% ----------------------------------------------------------
	% Introdução
	% ----------------------------------------------------------
	\chapter[Introdução]{Introdução}

Este Trabalho de Conclusão de Curso investiga a análise e a predição de trajetórias de espermatozoides em vídeos microscópicos por meio de técnicas de Visão Computacional. O estudo está inserido no contexto dos sistemas de análise seminal assistida por computador, conhecidos como \ac{CASA}, nos quais vídeos de amostras seminais são utilizados para observar o deslocamento celular e extrair informações relacionadas à motilidade. A automatização dessa análise exige transformar uma sequência de quadros em localizações, identidades persistentes e trajetórias mensuráveis.

Esse processo apresenta dificuldades próprias do domínio. As cabeças dos espermatozoides ocupam poucos pixels, têm aparência semelhante e podem apresentar baixo contraste em relação ao fundo. O movimento é rápido e irregular, e cruzamentos, sobreposições, detritos, variações de iluminação e deslocamentos globais da imagem podem produzir detecções incorretas ou trocas de identidade. O estudo utiliza exclusivamente vídeos microscópicos do VISEM e as anotações disponibilizadas pelo VISEM-Tracking para parte dessa coleção \cite{haugen2019visem,thambawita2023visemtracking}. O movimento observado será analisado em coordenadas de imagem, sem pressupor um escoamento imposto ou atribuir uma causa física à deriva. A proposta articula detecção, rastreamento multiobjeto, fluxo óptico e predição temporal em uma pipeline modular.

\section{Motivação}

A infertilidade constitui um problema de saúde pública com repercussões pessoais e sociais. A Organização Mundial da Saúde estima que aproximadamente uma em cada seis pessoas adultas vivencie infertilidade ao longo da vida \cite{who2023infertility}. Como a motilidade espermática integra a avaliação da fertilidade masculina, métodos computacionais que contribuam para uma análise mais objetiva, reprodutível e escalável podem apoiar estudos laboratoriais e pesquisas em reprodução humana. O presente trabalho, contudo, não pretende realizar diagnóstico clínico; seu objeto é a avaliação computacional de trajetórias em vídeos microscópicos.

Do ponto de vista acadêmico, o problema reúne desafios clássicos e atuais de Visão Computacional: detecção de objetos pequenos, associação temporal, manutenção de identificadores, estimação de movimento e modelagem de sequências. O conjunto VISEM reuniu vídeos microscópicos e dados de análise seminal de 85 participantes \cite{haugen2019visem}. Posteriormente, o VISEM-Tracking disponibilizou caixas delimitadoras, classes e identificadores persistentes em trechos de 30 segundos de 20 desses vídeos, viabilizando avaliações sistemáticas de detecção e \ac{MOT} \cite{thambawita2023visemtracking}. Esses trechos anotados permitem desenvolver e avaliar os módulos supervisionados. Os outros 65 vídeos da mesma coleção serão reservados à aplicação posterior e à análise qualitativa, pois não possuem anotações manuais de rastreamento para avaliar a acurácia das trajetórias.

Os trabalhos existentes demonstram avanços em partes do problema. O \textit{motilitAI} emprega fluxo óptico e aprendizado de máquina para estimar proporções agregadas de classes de motilidade \cite{ottl2022motilitai}; Noy et al. \cite{noy2023location} utilizam redes recorrentes para prever posições futuras a partir de trajetórias previamente obtidas; e Zhang et al. \cite{zhang2024spermyolo} combinam um detector especializado com rastreamento para o VISEM-Tracking. Permanece, entretanto, a oportunidade de avaliar de forma controlada se o movimento aparente estimado no entorno de cada célula acrescenta informação útil à predição individual, além do histórico de posições e velocidades.

\section{Problema}

O problema de pesquisa consiste em determinar se uma pipeline computacional consegue detectar e rastrear múltiplos espermatozoides em vídeos microscópicos e se características locais de fluxo óptico contribuem para prever suas posições futuras. A questão envolve dois níveis de dificuldade. No primeiro, erros do detector e do rastreador afetam a continuidade das trajetórias. No segundo, mesmo uma trajetória corretamente observada pode mudar de direção ou velocidade de modo que modelos cinemáticos simples deixem de representar sua dinâmica.

O fluxo óptico oferece uma descrição do deslocamento aparente de padrões de intensidade entre quadros, mas não corresponde automaticamente à velocidade física do fluido. Nos dados considerados, o campo observado pode misturar movimento das células, estruturas do meio, artefatos e movimento da câmera. Além disso, não há uma medição física de referência do escoamento. Assim, o trabalho precisa avaliar a utilidade computacional desse campo sem atribuir a ele uma interpretação física que os dados não sustentem.

A delimitação adotada concentra-se na localização bidimensional das cabeças dos espermatozoides e na predição de suas coordenadas em horizontes curtos. Morfologia detalhada, vitalidade, concentração seminal, diagnóstico médico e modelagem física completa do escoamento não fazem parte do escopo. Essa delimitação permite avaliar cada módulo com critérios objetivos e relacionar os resultados à hipótese de pesquisa.

\section{Hipótese}

A hipótese principal é que a inclusão de características locais de fluxo óptico, associadas ao histórico de posições de cada espermatozoide, reduz o erro de predição de posições futuras em comparação com modelos equivalentes que utilizam somente posição e velocidade. A verificação será feita por uma ablação pareada: serão mantidos os mesmos dados, divisões, horizontes, arquiteturas e hiperparâmetros, alterando-se apenas a presença das características de fluxo.

Três perguntas orientam a investigação:

\begin{enumerate}
	\item Qual combinação de detector e rastreador preserva melhor as identidades dos espermatozoides nas sequências avaliadas?
	\item O campo de movimento aparente apresenta consistência temporal suficiente para ser utilizado como atributo de predição?
	\item A inclusão desse campo reduz os erros de predição e em quais condições de densidade celular, duração de trajetória e horizonte temporal esse efeito ocorre?
\end{enumerate}

\section{Objetivos}

\subsection{Objetivo geral}

Investigar, implementar e avaliar uma pipeline computacional para análise e predição de trajetórias de espermatozoides em vídeos microscópicos, combinando técnicas de detecção, rastreamento multiobjeto e fluxo óptico, com ênfase na avaliação do efeito do movimento aparente sobre a previsão de posições futuras.

\subsection{Objetivos específicos}

\begin{enumerate}
	\item organizar os vídeos, metadados e anotações em divisões experimentais independentes e reproduzíveis;
	\item implementar e comparar abordagens clássicas e aprendidas para detectar espermatozoides;
	\item implementar e avaliar métodos de associação temporal capazes de manter identificadores persistentes;
	\item extrair trajetórias e calcular medidas cinemáticas compatíveis com a calibração disponível;
	\item estimar e analisar o campo de movimento aparente entre quadros, compensando movimento global quando necessário;
	\item comparar preditores de trajetórias com e sem características derivadas do fluxo óptico; e
	\item analisar quantitativa e qualitativamente os resultados, identificando limitações e condições de sucesso ou falha.
\end{enumerate}

\section{Contribuições esperadas}

A principal contribuição esperada é uma avaliação controlada da utilidade de características de fluxo óptico para a predição individual de trajetórias de espermatozoides. A comparação será feita tanto sobre trajetórias anotadas, isolando o preditor, quanto sobre trajetórias produzidas pela pipeline, permitindo medir a propagação dos erros de detecção e rastreamento.

Também se espera produzir uma organização reprodutível dos dados, linhas de base clássicas e modernas para cada módulo e uma análise das limitações envolvidas na transferência de métodos de Visão Computacional para vídeos microscópicos. O código, os modelos e os arquivos de resultados constituirão artefatos necessários aos experimentos, mas a contribuição científica será a evidência obtida sob critérios comuns de avaliação.

\section{Organização da Monografia}

Além desta introdução, a monografia está organizada em quatro capítulos. O Capítulo 2 apresenta os conceitos de análise seminal assistida por computador, processamento de imagens, detecção, rastreamento multiobjeto, fluxo óptico, predição de trajetórias, conjuntos de dados e métricas. O Capítulo 3 analisa os trabalhos relacionados e explicita a lacuna investigada. O Capítulo 4 descreve os dados, os módulos da pipeline, os experimentos e os critérios de validação. Por fim, o Capítulo 5 registra as alterações realizadas para consolidar as entregas anteriores e apresenta as expectativas para a continuidade do trabalho.

	
	
	% ----------------------------------------------------------
	% Fundamentação / Revisão Bibliográfica
	% ----------------------------------------------------------
	\chapter[Fundamentação Teórica]{Fundamentação Teórica}

Este capítulo apresenta os conceitos necessários para compreender a análise e a predição de trajetórias de espermatozoides em vídeos microscópicos. A organização acompanha o fluxo dos dados na pipeline: o domínio de aplicação define as medidas de interesse; o processamento de imagens permite localizar as células; o rastreamento associa observações ao longo do tempo; o fluxo óptico descreve movimento aparente; e a predição estima posições ainda não observadas. Ao final, são apresentados os conjuntos de dados e as métricas que permitem avaliar cada módulo separadamente.

\section{Análise seminal assistida por computador}

A análise seminal examina características do sêmen e dos espermatozoides, entre elas concentração, contagem total, morfologia, vitalidade e motilidade. A sexta edição do manual da Organização Mundial da Saúde sistematiza procedimentos de laboratório e ressalta a importância do controle de qualidade para produzir resultados comparáveis \cite{who2021manual}. Na avaliação básica, a motilidade pode ser classificada como progressiva, não progressiva ou imóvel.

Sistemas \ac{CASA} combinam microscopia, aquisição digital de vídeo e algoritmos para tornar essa avaliação quantitativa. Em vez de resumir a amostra somente pela inspeção visual, esses sistemas localizam cabeças de espermatozoides, associam observações da mesma célula e calculam grandezas cinemáticas a partir das coordenadas registradas no tempo. A automação, entretanto, não elimina a necessidade de padronização: taxa de quadros, ampliação, profundidade da câmara, contraste e regras de classificação influenciam diretamente os valores obtidos.

Considere uma trajetória bidimensional formada pelas posições $\mathbf{p}_t=(x_t,y_t)$, observadas em intervalos constantes $\Delta t$. Para uma trajetória com $T$ posições, a velocidade curvilínea, \ac{VCL}, corresponde ao comprimento acumulado do caminho dividido pelo tempo, enquanto a velocidade em linha reta, \ac{VSL}, considera apenas a distância entre a primeira e a última posição:

\begin{align}
	\operatorname{VCL} &=
	\frac{\sum_{t=1}^{T-1}\lVert \mathbf{p}_{t+1}-\mathbf{p}_t\rVert_2}
	{(T-1)\Delta t},\\
	\operatorname{VSL} &=
	\frac{\lVert \mathbf{p}_{T}-\mathbf{p}_1\rVert_2}
	{(T-1)\Delta t}.
\end{align}

A velocidade do percurso médio, \ac{VAP}, é calculada ao longo de uma versão suavizada da trajetória. A linearidade, \ac{LIN}, compara o deslocamento líquido com o caminho curvilíneo e pode ser expressa por $\operatorname{LIN}=100\,\operatorname{VSL}/\operatorname{VCL}$ \cite{who2021manual}. Para que as velocidades sejam apresentadas em $\mu$m/s, é necessário conhecer a escala espacial e a taxa de quadros. Sem essa calibração, as medidas devem permanecer em pixels por unidade de tempo.

\section{Processamento de imagens e detecção}

Uma imagem digital pode ser representada por uma função discreta de intensidade ou por múltiplos canais de cor. Em vídeos microscópicos, o processamento inicial busca aumentar a separação entre as cabeças dos espermatozoides e o fundo. Conversão para tons de cinza, redução de ruído, normalização de contraste, limiarização e operações morfológicas são transformações comuns nessa etapa \cite{szeliski2022computer}.

A segmentação por limiar divide os pixels de acordo com sua intensidade. O método de Otsu seleciona automaticamente um limiar que maximiza a separação estatística entre classes do histograma \cite{otsu1979threshold}. Após a binarização, erosão e dilatação podem remover componentes pequenos, preencher falhas ou separar estruturas segundo o elemento estruturante. Quando objetos próximos formam uma única região, a transformação \textit{watershed} oferece uma alternativa baseada em marcadores. Esses métodos são rápidos e interpretáveis, mas seus parâmetros são sensíveis a mudanças de contraste, iluminação, densidade celular e presença de detritos.

Na detecção de objetos, a saída contém, para cada instância, uma classe, uma caixa delimitadora e uma confiança. A família YOLO consolidou detectores de um estágio que realizam localização e classificação em uma única rede neural \cite{redmon2016yolo}. O VISEM-Tracking utiliza coordenadas normalizadas em formato compatível com YOLO e oferece classes para espermatozoides, aglomerados e células pequenas ou com cabeça reduzida \cite{thambawita2023visemtracking}.

A detecção de espermatozoides é um caso de detecção de objetos pequenos: poucos pixels descrevem cada cabeça, a textura individual é limitada e o borramento de movimento reduz a nitidez. Por isso, abordagens recentes utilizam resoluções maiores, camadas dedicadas a objetos pequenos e mecanismos de atenção. Zhang et al. \cite{zhang2024spermyolo}, por exemplo, adaptam o YOLOv8 e um rastreador para o domínio espermático. A avaliação precisa considerar precisão, revocação, qualidade de localização e tempo de execução, pois os erros do detector são propagados para o rastreamento.

\section{Rastreamento multiobjeto}

O \ac{MOT} consiste em estimar simultaneamente o estado de vários objetos e manter um identificador consistente para cada um ao longo do vídeo. Na estratégia \textit{tracking-by-detection}, cada quadro é processado por um detector e as novas observações são associadas às trajetórias existentes. O estado pode incluir posição, tamanho da caixa e velocidade. Uma trajetória é iniciada quando uma detecção persiste, atualizada enquanto recebe observações compatíveis e encerrada após permanecer ausente por um número definido de quadros.

Dois fundamentos desse processo são a estimação de estado e a associação de dados. O Filtro de Kalman representa a evolução do estado por um modelo linear com incerteza: primeiro prediz o estado no novo instante e depois corrige a previsão com a medição observada, ponderando as covariâncias dos ruídos de processo e de medição \cite{kalman1960filter}. Para associar trajetórias e detecções, constrói-se uma matriz de custos com distâncias entre centroides, sobreposição entre caixas ou outra medida de similaridade. O algoritmo Húngaro fornece uma atribuição global de custo mínimo sob correspondências um a um.

O SORT combina Filtro de Kalman, sobreposição entre caixas e associação Húngara em um rastreador simples e rápido \cite{bewley2016sort}. O ByteTrack, por sua vez, realiza associações sucessivas com detecções de alta e baixa confiança, buscando recuperar objetos que seriam descartados por um limiar único \cite{zhang2022bytetrack}. Esses métodos mostram a evolução de modelos baseados principalmente em movimento para estratégias que combinam localização, confiança e, em alguns casos, aparência.

No domínio espermático, a associação é dificultada pela alta densidade, pelas trajetórias rápidas e curvas e pela semelhança visual entre indivíduos. Descritores de reidentificação podem ser pouco discriminativos quando as cabeças ocupam poucos pixels; por outro lado, utilizar somente distância espacial falha em cruzamentos e mudanças bruscas de direção. As identidades manuais do VISEM-Tracking permitem medir trocas de identidade, fragmentações e recuperação após ausências \cite{thambawita2023visemtracking}.

\section{Fluxo óptico e movimento aparente}

Fluxo óptico é um campo vetorial que descreve o deslocamento aparente de padrões de intensidade entre quadros. Para cada posição $(x,y)$, procura-se um vetor $\mathbf{v}(x,y)=(u,v)$ que relacione a imagem no instante atual à imagem seguinte. Sob a hipótese de constância de brilho,

\begin{equation}
	I(x,y,t)=I(x+u\Delta t,y+v\Delta t,t+\Delta t),
\end{equation}

uma aproximação de primeira ordem produz a equação de restrição

\begin{equation}
	I_xu+I_yv+I_t=0.
\end{equation}

Há uma equação para duas incógnitas, portanto são necessárias restrições adicionais. Métodos locais, como Lucas--Kanade, supõem movimento aproximadamente constante em uma vizinhança. Formulações globais, como Horn--Schunck, introduzem suavidade no campo. Farnebäck aproxima vizinhanças por polinômios e estima fluxo denso em múltiplas escalas, constituindo uma linha de base prática para vídeos \cite{farneback2003twoframe,szeliski2022computer}.

Abordagens aprendidas estimam correspondências a partir de bases com deslocamento conhecido. O RAFT constrói volumes de correlação entre características dos dois quadros e atualiza iterativamente o campo por uma unidade recorrente \cite{teed2020raft}. Embora apresente bons resultados em conjuntos gerais de fluxo óptico, sua transferência para microscopia precisa ser avaliada, pois escala, textura, contraste e formação da imagem diferem das cenas de treinamento.

Neste trabalho, o fluxo possui dois usos potenciais: compensar movimento global da imagem e produzir características locais para a predição. É necessário distinguir esse campo da velocidade física do fluido. O fluxo óptico é medido em coordenadas de imagem e pode misturar deslocamento das células, do meio visível, de artefatos e da câmera. Como os dados do VISEM e as anotações do VISEM-Tracking não fornecem uma medição física de referência do escoamento, o resultado será denominado campo de movimento aparente.

\section{Predição de trajetórias}

Predição de trajetória consiste em estimar posições futuras a partir de uma sequência observada. Se $\mathbf{p}_{t-p+1:t}$ representa uma janela de $p$ posições, o objetivo é produzir $\widehat{\mathbf{p}}_{t+1:t+q}$ para um horizonte de $q$ quadros. A tarefa difere do rastreamento: o rastreador estima o estado atual com medições já disponíveis, enquanto o preditor é avaliado em posições ocultadas durante a inferência.

Modelos cinemáticos simples são referências importantes. A extrapolação de velocidade constante usa o deslocamento recente para projetar a posição, e o Filtro de Kalman acrescenta incertezas de processo e medição \cite{kalman1960filter}. Esses métodos são interpretáveis e exigem poucos dados, mas tendem a degradar quando a célula altera rapidamente sua direção ou velocidade.

Redes recorrentes representam dependências em sequências. A \ac{LSTM} utiliza células de memória e portas de controle para preservar informação por intervalos mais longos \cite{hochreiter1997lstm}. Em predição de trajetórias, posições ou deslocamentos sucessivos são codificados em um estado oculto, a partir do qual são geradas coordenadas futuras. O Social LSTM mostrou como essa formulação pode ser aplicada a trajetórias de agentes \cite{alahi2016sociallstm}; no domínio deste trabalho, Noy et al. \cite{noy2023location} aplicaram LSTM à localização futura de espermatozoides.

A hipótese computacional desta pesquisa considera que o histórico geométrico pode ser complementado por movimento aparente. Um modelo sem fluxo pode receber $(x,y,\Delta x,\Delta y)$, enquanto o modelo com fluxo acrescenta $(u,v,\lVert\mathbf{v}\rVert,\theta)$ amostrados na vizinhança da célula. A comparação deve manter arquitetura, histórico, horizonte e divisão dos dados controlados, alterando somente as características de fluxo.

\section{Conjuntos de dados}

O VISEM contém 85 vídeos microscópicos, um por participante, associados a resultados de análise seminal e atributos biológicos, mas sua publicação original não fornece identificadores de rastreamento quadro a quadro \cite{haugen2019visem}. O VISEM-Tracking acrescenta caixas delimitadoras, classes e identificadores persistentes verificados por especialistas a trechos iniciais de 30 segundos de 20 desses vídeos, totalizando 29.196 quadros \cite{thambawita2023visemtracking}. Trata-se de uma extensão anotada da mesma coleção de vídeos. Neste trabalho, os trechos anotados serão utilizados no desenvolvimento e na avaliação quantitativa de detecção, rastreamento e predição. Os demais 65 vídeos serão destinados à aplicação posterior e à análise qualitativa, sem atribuir acurácia às trajetórias na ausência de anotações manuais.

O protocolo de aquisição descreve amostras em platina aquecida a $37\,^{\circ}\mathrm{C}$, examinadas em microscópio Olympus CX31 com aumento de $400\times$ e registradas por câmera UEye UI-2210C acoplada ao microscópio \cite{haugen2019visem,thambawita2023visemtracking}. Essas publicações não documentam tubo, bombeamento ou contracorrente induzida durante a filmagem. O artigo original relata que, para alguns participantes, foram realizadas duas gravações devido à deriva na primeira amostra registrada, mantendo-se um vídeo por participante na coleção \cite{haugen2019visem}. Esse relato não estabelece a origem física da deriva nem comprova um escoamento controlado; a ausência de tal descrição também não demonstra que o líquido esteja imóvel.

As divisões de treinamento, validação e teste serão feitas por vídeo de origem, mantendo os trechos do mesmo participante na mesma divisão, nunca por quadros aleatórios. Clipes derivados não serão tratados como amostras independentes. Essa precaução impede que imagens quase idênticas da mesma sequência apareçam simultaneamente no treinamento e no teste.

\section{Métricas de avaliação}

Na detecção, precisão e revocação são definidas a partir de verdadeiros positivos (VP), falsos positivos (FP) e falsos negativos (FN):

\begin{equation}
	\text{Precisão}=\frac{\mathrm{VP}}{\mathrm{VP}+\mathrm{FP}},
	\qquad
	\text{Revocação}=\frac{\mathrm{VP}}{\mathrm{VP}+\mathrm{FN}}.
\end{equation}

Uma detecção é considerada correspondente ao gabarito quando a \ac{IoU} ultrapassa um limiar. A área sob a curva de precisão--revocação gera a precisão média, e a média entre classes ou limiares gera a \ac{mAP}. Como poucos pixels de deslocamento afetam significativamente a IoU de objetos pequenos, também é útil analisar erro de centroide e resultados por classe.

No rastreamento, a \ac{MOTA} combina falsos positivos, falsos negativos e trocas de identidade. A \ac{IDF1} mede a preservação de identidade por meio da média harmônica entre precisão e revocação no espaço das associações. A \ac{HOTA} foi proposta para equilibrar qualidade de detecção, associação e localização \cite{luiten2021hota}. Neste trabalho, HOTA será a métrica principal de rastreamento, acompanhada por MOTA, IDF1, trocas de identidade, fragmentações e tempo por quadro.

Na predição, o \ac{ADE} resume o erro ao longo de todo o horizonte, enquanto o \ac{FDE} considera a última posição prevista. Para $N$ trajetórias e horizonte de $q$ quadros:

\begin{align}
	\operatorname{ADE} &=
	\frac{1}{Nq}\sum_{i=1}^{N}\sum_{k=1}^{q}
	\left\lVert\widehat{\mathbf{p}}_{i,t+k}-\mathbf{p}_{i,t+k}\right\rVert_2,\\
	\operatorname{FDE} &=
	\frac{1}{N}\sum_{i=1}^{N}
	\left\lVert\widehat{\mathbf{p}}_{i,t+q}-\mathbf{p}_{i,t+q}\right\rVert_2.
\end{align}

Essas medidas devem ser apresentadas no mesmo sistema de coordenadas e, quando houver calibração, também em unidade física \cite{alahi2016sociallstm}. Médias serão acompanhadas por dispersão, intervalos de confiança e análise estratificada por horizonte, densidade celular e duração da trajetória.


	% ----------------------------------------------------------
	% Trabalhos relacionados
	% ----------------------------------------------------------
	\chapter[Trabalhos Relacionados]{Trabalhos Relacionados}

Este capítulo analisa quatro trabalhos diretamente relacionados às etapas da pipeline proposta. Os estudos foram selecionados por abordarem motilidade espermática, disponibilização de trajetórias anotadas, predição de posições ou detecção e rastreamento no mesmo domínio. A comparação considera objetivo, dados, método, resultados e aspectos ainda não tratados em conjunto.

Ottl et al. \cite{ottl2022motilitai} propõem o \textit{motilitAI}, um arcabouço para estimar automaticamente as proporções de espermatozoides progressivos, não progressivos e imóveis em vídeos microscópicos. O método realiza pré-processamento dos vídeos do VISEM, rastreia células por uma adaptação esparsa do fluxo óptico de Lucas--Kanade, extrai descritores de movimento e aparência e utiliza modelos de regressão. A melhor configuração alcançou erro absoluto médio de 7,31 pontos percentuais, inferior aos 8,83 pontos do melhor resultado da competição de referência utilizada pelos autores. O estudo demonstra que movimento e trajetórias extraídos de vídeo fornecem informação sobre motilidade, porém sua saída é agregada para a amostra. O presente projeto, em contraste, pretende manter trajetórias individuais e prever posições futuras.

Thambawita et al. \cite{thambawita2023visemtracking} apresentam o VISEM-Tracking, criado para reduzir a escassez de vídeos públicos com anotações temporais de espermatozoides humanos. O conjunto contém 20 vídeos de 30 segundos, totalizando 29.196 quadros, com caixas delimitadoras, identificadores persistentes e classes para espermatozoides, aglomerados e células pequenas ou com cabeça reduzida. Os autores também disponibilizam características seminais e sequências sem rótulos. Como validação técnica, foram treinadas variantes do YOLOv5, demonstrando que as anotações sustentam modelos supervisionados. A relação com este projeto é direta, pois os dados podem apoiar os módulos de detecção, rastreamento e predição. O artigo, contudo, não propõe um preditor nem estima movimento do meio.

Noy et al. \cite{noy2023location} investigam a localização futura de espermatozoides para aplicações que exigem posicionamento rápido e não invasivo. Os autores combinam uma rede LSTM multicamadas, responsável por aprender a dinâmica não linear a partir do histórico de posições, com extrapolação linear e uma função de perda ajustada ao problema. Os experimentos utilizam 501 trajetórias, com 114.229 quadros obtidos de três doadores a 70 quadros por segundo, e avaliam horizontes de até 1,43 segundo. Para horizontes inferiores a 0,6 segundo, o erro médio de localização permanece entre 4 e 8~$\mu$m; em várias configurações, o modelo reduz o erro da extrapolação linear. O estudo constitui uma referência direta para a etapa de predição, mas parte de trajetórias previamente obtidas e não incorpora detecção multiobjeto ou características de fluxo óptico.

Zhang et al. \cite{zhang2024spermyolo} tratam conjuntamente a detecção e o rastreamento de espermatozoides com o método Sperm YOLOv8E-TrackEVD. O detector adapta o YOLOv8 com mecanismo de atenção, camada dedicada a objetos pequenos e módulos especializados. O rastreador modifica o DeepOCSORT para melhorar a associação temporal nesse domínio. No VISEM-Tracking, o sistema completo alcança HOTA de 74,303\% e MOTA de 71,167\%. O trabalho fornece uma referência para os módulos de detecção, associação e avaliação, mas não aborda a estimação de movimento aparente do meio nem a predição de posições futuras.

\section{Comparação entre os trabalhos}

A Tabela \ref{tab:comparacao-trabalhos} sintetiza as diferenças entre os estudos. Em conjunto, eles fornecem dados, métodos e métricas para as partes centrais da proposta, mas nenhum avalia, em uma mesma pipeline, a influência de características locais de fluxo óptico sobre a predição individual após a detecção e o rastreamento.

\begin{table}[H]
\centering
\caption{Comparação dos trabalhos relacionados}
\label{tab:comparacao-trabalhos}
\footnotesize
\renewcommand{\arraystretch}{1.25}
\begin{tabularx}{\textwidth}{@{}p{2.6cm}YYY@{}}
\toprule
\textbf{Trabalho} & \textbf{Objetivo e dados} & \textbf{Método e resultado principal} & \textbf{Relação com este projeto} \\
\midrule
Ottl et al. (2022) &
Estimar proporções de classes de motilidade em vídeos do VISEM &
Fluxo óptico esparso, descritores e regressão; erro absoluto médio de 7,31 pontos percentuais &
Utiliza movimento e rastreamento, mas produz apenas estimativas agregadas da amostra \\
\addlinespace
Thambawita et al. (2023) &
Disponibilizar vídeos, caixas e identidades de rastreamento &
Anotação especializada de 20 vídeos e linhas de base com YOLOv5 &
Fornece dados para detecção, rastreamento e avaliação de trajetórias \\
\addlinespace
Noy et al. (2023) &
Prever posições futuras em 501 trajetórias &
LSTM combinada à extrapolação linear; erro de 4 a 8~$\mu$m em horizontes menores que 0,6 s &
É referência de predição, mas não inclui detecção, MOT ou características de fluxo \\
\addlinespace
Zhang et al. (2024) &
Detectar e rastrear no VISEM-Tracking &
YOLOv8 aprimorado e DeepOCSORT modificado; HOTA de 74,303\% e MOTA de 71,167\% &
É referência para detecção e associação, sem prever trajetórias futuras \\
\bottomrule
\end{tabularx}
\fonte{Elaborada pelo autor (2026).}
\end{table}

\section{Lacuna investigada}

Os estudos analisados evidenciam três lacunas complementares. Primeiro, trabalhos de análise de motilidade podem utilizar trajetórias e fluxo óptico sem preservar uma previsão individual de cada célula. Segundo, conjuntos e métodos de rastreamento fornecem identidades persistentes, mas encerram a análise no estado observado. Terceiro, modelos de predição partem de trajetórias prontas e não medem como erros da pipeline afetam o horizonte futuro.

O presente trabalho procura conectar essas etapas por meio de uma arquitetura modular. A contribuição não consiste apenas em encadear algoritmos, mas em comparar, sob os mesmos dados e critérios, preditores com e sem características de movimento aparente. A avaliação separada sobre trajetórias anotadas e trajetórias estimadas permitirá distinguir a capacidade do preditor da degradação introduzida por detecção e rastreamento.

		
	% ----------------------------------------------------------
	% Método
	% ----------------------------------------------------------
	\chapter[Método]{Método}
\label{cap:metodo}

O trabalho será conduzido como uma pesquisa aplicada, quantitativa e experimental em Visão Computacional. O método consiste na implementação e na avaliação controlada de uma pipeline para detecção, rastreamento multiobjeto, estimação de movimento aparente e predição de trajetórias de espermatozoides em vídeos microscópicos. A hipótese será examinada por experimentos reproduzíveis que comparem modelos equivalentes com e sem características locais de fluxo óptico.

\section{Produtos e dados da pesquisa}

Os produtos derivados serão uma coleção de dados organizada com divisões experimentais documentadas; uma pipeline modular em Python; modelos treinados; arquivos contendo caixas, identificadores, trajetórias e campos vetoriais; e uma análise comparativa apresentada por métricas, tabelas e visualizações. O código e os modelos serão artefatos necessários à realização dos experimentos, enquanto o produto científico principal será a evidência sobre o efeito do fluxo óptico na predição de trajetórias.

Os vídeos microscópicos utilizados serão exclusivamente os do VISEM, que contém 85 vídeos, um por participante \cite{haugen2019visem}. O VISEM-Tracking fornece caixas, classes e identificadores persistentes por quadro para trechos iniciais de 30 segundos de 20 desses vídeos \cite{thambawita2023visemtracking}. Os trechos anotados dessas 20 unidades constituirão o gabarito para desenvolvimento e avaliação quantitativa. Os outros 65 vídeos serão processados somente após o congelamento das configurações, como aplicação e análise qualitativa, pois não possuem tracking manual que permita declarar acurácia. Os 502 clipes sem rótulos extraídos dos mesmos vídeos serão tratados como derivados, e não como amostras independentes. Intervalos sem arquivo de anotação no vídeo 23 serão excluídos de treino e métricas, nunca interpretados como quadros negativos. A avaliação não pressupõe contracorrente induzida nem inclui uma segunda coleção experimental com tubo.

\section{Levantamento, geração e organização dos dados}

Para cada vídeo serão registrados resolução, taxa de quadros, duração, escala espacial quando disponível e condições de aquisição. Os quadros serão extraídos preservando a ordem temporal e vinculados a um identificador único. O pré-processamento será limitado a operações cuja influência possa ser controlada, como conversão de cor, normalização de intensidade, correção moderada de contraste e redução de ruído. As anotações serão convertidas para uma representação comum, e as coordenadas das trajetórias serão obtidas a partir do centro das caixas.

Será adotada uma separação fixa por vídeo em 12 unidades de treinamento, quatro de validação e quatro de teste. O treinamento servirá para estimar parâmetros; a validação, para escolher hiperparâmetros e critérios de parada; e a bateria confirmatória no teste ocorrerá somente após o congelamento da configuração. A confirmação em cinco folds dependerá de um desenho explícito que separe ajuste e seleção dos vídeos avaliados em cada fold. O piloto anterior examinou frames dos quatro vídeos hoje reservados ao teste; portanto, esse holdout não é completamente cego. Repetir uma configuração escolhida nos cinco grupos ou redistribuir os vídeos não restaura a independência perdida por essa exposição histórica, que será declarada como limitação. Arquivos de configuração registrarão divisões, sementes 42, 123 e 2026, versões das bibliotecas, hardware e parâmetros resolvidos.

\section{Detecção e rastreamento multiobjeto}

A detecção começará com baselines clássicos separados: threshold fixo com morfologia, Otsu, threshold adaptativo, Blob Detector, MOG2, KNN e Watershed. Uma variante híbrida com correção de fundo e normalização local será avaliada sem substituir os clássicos puros. Por fim, será treinado um detector da família YOLO com as anotações do VISEM-Tracking. A métrica principal será o F1 de localização dos indivíduos anotados nas classes 0 e 2, obtido por matching Húngaro um-para-um entre centros, com tolerância de 10 pixels na resolução original. Todas as previsões participam da associação, independentemente da classe prevista. Após a associação, previsões residuais a até esse raio de qualquer indivíduo anotado permanecem falsos positivos; das restantes, são ignoradas somente aquelas cujo centro está dentro de uma caixa de agrupamento da classe 1. Os indivíduos anotados dentro dessas caixas continuam sendo avaliados. A regra será recalculada integralmente nas sensibilidades de 15 e 20 pixels. Uma avaliação complementar considerará todos os objetos anotados, incluindo cada agrupamento como um objeto, sem interpretá-lo como uma única célula. Serão informadas as contagens brutas, avaliadas e ignoradas, a cobertura espacial das regiões de agrupamento, precisão, revocação, erro de centro, erro de contagem, latência e memória. A MAE de contagem utilizará o mesmo universo de quadros anotados para todos os detectores. As caixas manuais de agrupamento serão usadas exclusivamente na avaliação, sem filtrar entradas ou saídas do detector e do rastreador. As coordenadas, dimensões e valores de confiança serão exportados com a precisão de ponto flutuante usada no processamento, permitindo reconstruir as associações e os erros a partir dos arquivos de saída; arredondamentos serão restritos à apresentação. A tolerância é uma convenção operacional de localização em relação ao centro geométrico da caixa anotada, não uma estimativa de incerteza dos anotadores ou do centro anatômico exato. mAP@0,5, mAP@0,5:0,95 e a análise nas três classes serão secundários e específicos do YOLO. O método de Zhang et al. \cite{zhang2024spermyolo} será utilizado como referência do domínio, considerando as diferenças de dados e ambiente computacional.

Para a triagem prospectiva do threshold fixo, foi registrado um espaço de 171 configurações: limiares de 0 a 255 em passos de 16, incluindo adicionalmente 190 e 200 como referências históricas, com abertura e fechamento de zero, uma ou duas iterações. O limiar 255 encerra a faixa; a sequência de passo 16 termina em 240. Serão mantidos kernel elíptico de $3\times3$, ausência de suavização adicional e filtro de área de 3 a 300 pixels. Uma amostra de 48 quadros anotados por vídeo de treino será selecionada por posições igualmente espaçadas no conjunto ordenado dos índices elegíveis; 12 desses quadros por vídeo formarão o subconjunto comum da triagem, e um quadro desse subconjunto será reservado à calibração de custo. Os pixels serão obtidos por decodificação sequencial dos MP4 canônicos, sem redimensionamento ou nova compressão, com hashes das entradas e das anotações.

Todas as configurações serão comparadas nos mesmos 144 quadros, usando a média dos F1 agregados dentro de cada um dos 12 vídeos. Desempates considerarão revocação, menor MAE de contagem e identificador, nessa ordem; tempo de execução não participa da escolha. Os cinco primeiros candidatos de uma bateria completa seguirão para um refinamento com limiares inteiros entre menos 15 e mais 15 em torno de cada candidato, mantendo sua morfologia e usando os 48 quadros por vídeo. Configurações repetidas serão deduplicadas, com limite de 155 candidatos. As duas finalistas desse refinamento poderão ser comparadas nos vídeos completos de validação. Essas etapas constituem seleção no treino e na validação; seus resultados não estimam desempenho em dados independentes. O protocolo registra orçamento e critérios antes da execução e interrompe comparações incompletas, sem classificar somente os candidatos que terminaram.

Antes de executar o refinamento, será realizado um novo benchmark com todos os candidatos nos mesmos 12 quadros reservados à medição de custo. A projeção somará o tempo de validação das entradas a 96 vezes o tempo do laço desse benchmark: fator 48 pela ampliação da amostra e fator 2 como margem operacional. O limite de liberação será de 4.800 segundos, com teto de 7.200 segundos para a bateria de refinamento. A projeção não constitui garantia de pior caso temporal. Todos os candidatos serão reavaliados nos 576 quadros; não serão somadas métricas antigas a resultados parciais novos. A escolha das duas finalistas seguirá somente a ordenação registrada, sem exigir diversidade de morfologia. Como a amostra passa de 12 para 48 quadros por vídeo, diferenças entre os valores agregados dessas duas etapas não serão interpretadas isoladamente como melhora do detector.

Em 8 de setembro de 2026, a triagem e o refinamento do threshold v3 estavam concluídos e conferidos. O refinamento avaliou 117 configurações nos mesmos 576 quadros de treino, resultando em T219/o0/c2 e T218/o0/c2 como finalistas, sem promoção do detector ou conclusão sobre a hipótese de fluxo. O protocolo da bateria subsequente previa comparar somente essas duas configurações nos quatro vídeos completos de validação: 14, 19 e 36, com 1.470 quadros cada, e 52, com 1.440, totalizando 5.850 quadros físicos e 11.700 avaliações de configuração--quadro. A seleção prevista seguia F1 macro por vídeo a 10 pixels, revocação, menor MAE de contagem e identificador; sensibilidades de 15/20 pixels e custo tinham caráter descritivo. Não estavam previstos testes de hipótese ou intervalos de confiança nessa etapa de seleção, nem congelamento automático ou acesso ao teste. A conclusão posterior dessa bateria é registrada abaixo.

O contrato anterior a essa execução exigia código e protocolo registrados em commit, com controles testados para exigir a quantidade exata de quadros e detectar quadros extras, rejeitar anotação ausente ou malformada e conferir hashes das entradas e das configurações. Esses controles ainda estavam em correção na revisão anterior à bateria. O orçamento prospectivo então fixado era de 1.800 segundos de tempo de parede, com controle operacional suave, limites de 2 GiB de RSS e de artefatos, e 2.000 previsões por quadro. A ordem dos oito pares configuração--vídeo seria embaralhada com seed 42, sem interpretar a ordem como repetição independente, e não seriam gerados vídeos de saída. Falhas deveriam preservar os artefatos parciais e impedir a seleção até a conclusão e conferência de todos os pares.

As detecções serão associadas temporalmente por métodos de complexidade crescente: centroide guloso, associação Húngara global, SORT, que combina Filtro de Kalman e associação Húngara \cite{bewley2016sort}, e uma implementação ByteTrack-style que reaproveita detecções de baixa confiança \cite{zhang2022bytetrack}. Uma variante Adaptive Flow-SORT será tratada separadamente como híbrida. Cada tracker será avaliado primeiro com caixas do ground truth, isolando a associação, e depois com o detector congelado, medindo o cenário fim a fim. Limiares de associação e o número de quadros tolerados sem detecção serão escolhidos na validação. As saídas terão a forma $(v,t,i,x,y)$, em que $v$ identifica o vídeo, $t$ o quadro, $i$ a célula e $(x,y)$ a posição de seu centro.

Antes da avaliação de identidades e da construção de janelas, será especificada a elegibilidade das trajetórias individuais das classes 0 e 2, com tratamento explícito das anotações de agrupamento, mudanças de classe e lacunas. Uma trajetória de caixa de agrupamento não será apresentada como trajetória de uma única célula. O protocolo de associação temporal e a preparação das entradas para o HOTA oficial serão documentados separadamente, sem transferir automaticamente a tolerância de 10 pixels usada na detecção.

Em seguida, a validação completa foi executada sob o commit \texttt{7f47afb}, com Git limpo e 824 testes aprovados antes da bateria. As oito runs somaram 11.700 avaliações em 166,657535 segundos, com pico amostrado de RSS de 236,594 MiB. T218/o0/c2 foi selecionada pelo critério prospectivo, com F1 macro de 0,658959, revocação macro de 0,794711 e MAE de contagem macro de 4,408433. T219/o0/c2 obteve F1 macro de 0,657819. A diferença de 0,11403 ponto percentual é descritiva: T218 teve F1 maior em 19/36, e T219 em 14/52. A conferência independente reconstruiu os 11.700 registros e agregados, com 1.772.878 comparações de campos/valores e 144 verificações amostrais de matching com SciPy. Duas tentativas iniciais foram preservadas após correções de suposições do verificador sobre formato de metadados; nenhuma run científica foi alterada. Na conclusão da bateria, a seleção ainda não estava congelada e não constituía teste independente, avaliação de folds ou evidência sobre a contribuição do fluxo. Posteriormente, T218/o0/c2 foi fixada como linha de base de desenvolvimento, com parâmetros e avaliação preservados. Esse escopo mantém teste, folds e aplicação bloqueados até um desenho confirmatório próprio, sem restaurar a cegueira historicamente perdida.

\subsection{Busca entre famílias de detectores: 11 de setembro de 2026}
\label{sec:busca-classicos-20260911}

A comparação entre famílias foi iniciada por uma busca registrada no treino, sob o commit \texttt{2547109}. Foram avaliadas 43 configurações em seis famílias: threshold fixo T218/o0/c2 como referência de desenvolvimento, Otsu, threshold adaptativo gaussiano, realce com CLAHE seguido de threshold, Blob e Watershed. Todas receberam os mesmos 576 quadros anotados, 48 por vídeo nos 12 treinos, reutilizados do cache autenticado, totalizando 24.768 avaliações de configuração--quadro. A métrica principal foi a média dos F1 por vídeo a 10 pixels, após agregação das contagens dentro de cada vídeo, com sensibilidades de 15/20 pixels e avaliação secundária de todos os objetos. Quadros, configurações e observações repetidas não foram tratados como unidades independentes.

O primeiro smoke foi interrompido pelo teto de 2.000 previsões por quadro, sem seleção. O plano e a execução parcial foram preservados. Uma revisão operacional explícita elevou somente esse teto para 307.200 previsões por quadro, mantendo a grade científica, os dados, as métricas e os limites de 2 GiB de memória e de artefatos. Um novo smoke completo, com 516 avaliações, precedeu a busca. Essa revisão não alterou as regras da execução histórica nem truncou previsões para concluir a comparação.

Na grade avaliada, \texttt{blob\_v1\_003} obteve F1 macro de 0,791005 e T218/o0/c2 obteve 0,775634, diferença de aproximadamente 1,54 ponto percentual. Blob apresentou F1 maior em sete dos 12 vídeos, e T218 nos outros cinco. Portanto, a diferença média não representa superioridade uniforme. As famílias tiveram quantidades diferentes de configurações pesquisadas, e T218 já havia sido ajustado em uma busca anterior. O resultado descreve esta triagem limitada no treino, não o melhor desempenho possível de cada família ou sua generalização. Foram retidas duas finalistas por família pesquisada e a referência T218, totalizando 11 candidatas à próxima etapa, sem promoção, liberação de validação ou acesso ao teste.

A busca terminou em 752,160737 segundos, com RSS máximo amostrado de 823,145 MiB. A conferência independente verificou 289 arquivos, 30.927.965 comparações e 148.608 associações com SciPy, sem diferença numérica nos valores conferidos. A paridade de T218 nos mesmos quadros foi confirmada contra o refinamento anterior. Essa conferência autentica e reconstrói os derivados; não reexecuta os detectores nem a decodificação original. A comparação ainda não inclui YOLO, MOG2 ou KNN: o primeiro exige treinamento próprio, e os métodos temporais precisam de processamento contínuo e aquecimento, incompatíveis com tratar os quadros espaçados como imagens independentes. O refinamento e a validação das alternativas exigirão protocolos novos; maior F1 de detecção, por si só, não assegura maior HOTA ou menor erro de predição.


\subsection{Refinamento local e preparação dos dados aprendidos}
\label{sec:refinamento-classicos-20260911}

Após a busca anterior, foram registrados o contrato e a implementação de um refinamento local, antes das novas inferências. As duas finalistas de cada uma das cinco famílias pesquisadas originaram 54 propostas que incluíam os próprios pais e alterações de um parâmetro por vez. Três propostas ficaram fora dos limites prospectivos; as 51 válidas resultaram em 45 configurações únicas após deduplicação, preservando todas as origens. A referência T218 foi autenticada e não reexecutada nesta etapa.

Um smoke de 540 avaliações precedeu o refinamento completo. Ambos foram concluídos e conferidos independentemente sob o commit \texttt{ca68f16}. O refinamento utilizou os mesmos 576 quadros físicos dos 12 treinos, totalizando 25.920 avaliações. A melhor configuração da grade de cada família apresentou os seguintes F1 macro a 10 pixels: Blob: 0,791005; Watershed: 0,627005; Otsu: 0,626315; Adaptativo: 0,446885; Híbrido CLAHE: 0,364491. Essas médias descrevem seleção nos dados de treino, com esforço desigual entre famílias. Foram retidas duas finalistas por família, totalizando dez, sem promoção, abertura do teste ou escolha final da pipeline.

Os dez pais reproduziram suas detecções, anotações exportadas e métricas científicas nos mesmos quadros, exceto tempo e identificação da nova execução. A conferência independente reconstruiu vizinhança, associações, agregações e classificação. Esse controle verifica consistência e reprodutibilidade dos derivados; não repete a inferência nem certifica generalização. A duração da bateria foi 528,244100 segundos e o máximo de RSS amostrado foi 410,406000 MiB. Recursos e custos permanecem vinculados ao ambiente e ao commit da execução.

Foi também preparado um dataset derivado independente para o futuro treinamento YOLO: 17.466 pares de imagem e anotação no treino e 5.850 na validação, totalizando 23.316. As 174 lacunas do vídeo 23 foram excluídas, sem tratá-las como negativos. As três classes foram preservadas; os rótulos geométricos YOLO foram confrontados com as anotações FTID usadas na avaliação, ignorando somente o identificador. Fontes, cópias e referências foram autenticadas por hashes, e as cópias JPEG foram decodificadas e conferidas independentemente. Não houve leitura de imagens do teste ou treinamento de rede nesta preparação.

O dataset preparado permanece imutável. O futuro executor de treinamento deverá usar outra cópia independente para isolar os caches e eventuais reparos da biblioteca, mantendo rastreável qualquer mudança de entrada. Receita, pesos iniciais, sementes, seleção de checkpoints e parâmetros de inferência exigem contrato próprio. A sequência seguinte é validar as finalistas clássicas em vídeos completos e treinar e comparar YOLO sob a mesma regra de detecção; os métodos temporais requerem protocolo contínuo com aquecimento.

A paridade geométrica dos rótulos não certifica igualdade de pixels entre os JPEGs recebidos e os quadros decodificados dos vídeos. Os JPEGs constituem as entradas de treinamento preparadas; a comparação de detecção entre YOLO e os clássicos deverá processar os mesmos quadros do vídeo ou cache de avaliação, com pré-processamento explicitamente registrado.

\section{Estimação do movimento aparente}

O movimento aparente será estimado por Lucas--Kanade, Farnebäck \cite{farneback2003twoframe}, Horn--Schunck e, se os recursos computacionais permitirem, RAFT \cite{teed2020raft}. Variantes híbridas combinarão compensação de movimento global, consistência direto--reverso, rejeição robusta e eventual refinamento neural. Regiões ocupadas pelas células e por artefatos evidentes serão mascaradas quando o objetivo for caracterizar o fundo; o clássico puro e cada adaptação manterão configuração e resultado próprios.

Como os dados do VISEM e as anotações do VISEM-Tracking não fornecem uma medição física do campo de velocidade do fluido, o resultado será tratado como campo de movimento aparente. Sua consistência será analisada por erro fotométrico após a deformação de um quadro em direção ao seguinte, concordância entre fluxos direto e reverso, regularidade temporal e estabilidade nas regiões de fundo. Sequências sintéticas com deslocamentos conhecidos poderão complementar a verificação controlada dos estimadores, sem constituir uma segunda coleção de vídeos microscópicos reais. A conversão dos deslocamentos aparentes para unidades físicas somente será realizada quando houver calibração espacial e temporal confiável; essa conversão, por si só, não estabelece a velocidade física do fluido.


Como verificação inicial da conexão causal com a referência individual, foi
registrado e executado um ensaio de Farnebäck sob o commit \texttt{16eecbb},
após 1.495 testes automatizados. Foram utilizados somente os quadros de 0 a 19
dos vídeos de treino 11 e 12, sem máscara ou busca de parâmetros. Na origem
$t=19$, cada um dos 19 campos históricos $F_s$, que representa $s\to s+1$,
foi amostrado no centro anotado $\mathbf p_s$ do primeiro quadro do par, com
interpolação bilinear e pesos em precisão de 64 bits. O último par disponível
é $18\to19$; nenhuma imagem posterior foi fornecida ao estimador. O consumidor
recebeu apenas o histórico e as chaves das janelas, sem coordenadas futuras.
A elegibilidade continua condicionada à disponibilidade de futuro completo
na referência derivada, exclusivamente para avaliação retrospectiva.

As 68 janelas selecionadas antes dos pixels produziram 1.292 amostras com
suporte numérico válido, sem exclusões. Os 38 pares foram estimados nos dois
sentidos, totalizando 76 campos, além de 36 comparações temporais. Os erros
fotométricos médios após alinhamento foram 0,002374 e 0,003461 nos vídeos 11
e 12, contra 0,003008 e 0,005554 com deslocamento zero, em intensidade cinza
dividida por 255 e sobre o mesmo suporte de cada par. As consistências
direto--reverso médias foram 0,015067 e 0,169551 pixels. As médias por vídeo
dão peso igual aos pares com suporte. Validade numérica de todas as amostras
não implica acurácia de 100\%, e esses diagnósticos não medem a velocidade
física do fluido nem o ganho em predição.

O ensaio levou 56,923953 segundos, com RSS máximo amostrado de 289,598 MiB.
A conferência independente verificou 190 arquivos e 37.406 comparações,
das quais 8.378 numéricas, com diferença máxima de $8{,}88\times10^{-16}$,
abaixo da tolerância absoluta de $10^{-9}$. Uma falha inicial na leitura do
esquema do relatório histórico da referência foi preservada e corrigida sob
\texttt{0fccda8}, após nove testes sintéticos, sem modificar ou repetir a run.
A conferência reconstituiu a matemática dos derivados, sem importar a
implementação científica; não constitui decodificação independente do MP4 ou
reexecução do estimador. O ensaio certifica a conexão operacional nesse
prefixo de dois vídeos. A extração ampliada e a ablação pareada com e sem fluxo
permanecem sujeitas a contrato prospectivo próprio.


Para preparar a escala, foi registrado antes dos pixels um benchmark de
armazenamento compacto sob o commit \texttt{6110c44}, após 1.728 testes.
Foram processados os quadros de 0 a 59 de cada um dos 12 vídeos de treino,
com a mesma configuração de Farnebäck, uma thread de CPU e OpenCL desativado.
Os 708 campos diretos foram amostrados nas posições históricas de origem.
Foram retidos campos densos nos dois sentidos apenas nos pares $0\to1$ e
$58\to59$ de cada vídeo, além de seus quadros cinza. As 10.848 janelas
compartilham 15.762 amostras distintas, correspondentes a 206.112 usos.
Cada amostra preserva os quatro vizinhos usados na interpolação; o QA pode
reconstruir seus pesos e valores, mas somente nos sentinelas pode confrontar
os vizinhos com o campo denso retido. Não houve segunda decodificação ou
reexecução independente do estimador. Todas as amostras tiveram suporte
numérico válido, sem que isso certifique acurácia do movimento estimado.

O benchmark terminou em 199,583442 segundos, com RSS amostrado de 375,145 MiB
e 154.023.508 bytes finais. A conferência independente passou na primeira
tentativa: 288 arquivos e 934.668 comparações, com diferença máxima de
$3{,}55\times10^{-15}$. A regra de projeção, fixada antes dos resultados e
com fator de segurança dois, estimou 8.052,675671 segundos e 1.768,62 MiB para
a escala completa. O tempo excedeu o teto prospectivo de 7.200 segundos;
o armazenamento permaneceu abaixo dos 4.096 MiB. Por isso, a extração completa
não foi liberada por aquele protocolo. Na conclusão do benchmark, a continuidade
prevista era discriminar custos de decodificação, estimação e amostragem antes
de ampliar a extração. Em 11 de setembro, o tempo deixou de ser impedimento
para novos planos: o teto e a projeção históricos permanecem preservados, mas
uma nova execução pode monitorar tempo sem aquele corte, mantendo controle de
memória, armazenamento e completude. A instrumentação deixa de ser condição
obrigatória para a ablação exploratória nas 10.848 janelas já disponíveis,
que ainda exige contrato e executor próprios. Não houve avaliação de ADE/FDE
com fluxo, seleção de parâmetros ou teste da hipótese neste benchmark.

\section{Predição das trajetórias}

Antes dos preditores, foi registrado e executado o protocolo de referência individual apenas nos 12 vídeos de treino, sob o commit \texttt{33d191d}, com Git limpo e 1.026 testes aprovados. As classes 0 e 2 preservam o ID original; classe 1, ausência do ID e ausência de anotação interrompem os segmentos. Mudanças entre classes 0 e 2 não interrompem o trecho. A segmentação não renumera identidades biológicas e não define implicitamente a política de HOTA. Todas as observações e todos os segmentos são mantidos, inclusive os curtos.

Foram preservadas 363.074 observações individuais em 725 segmentos, dos quais 623 geram 343.776 janelas de 20 posições observadas e dez posições futuras, com passo de um quadro. Os 102 segmentos sem janela conservam 1.231 observações. Das origens individuais, 12.914 não têm histórico completo e 6.384 têm histórico, mas não o futuro completo. A elegibilidade é condicionada à disponibilidade de referência individual futura; esse uso é exclusivo dos alvos e da avaliação retrospectiva, não das entradas do preditor. As janelas permanecem dependentes dentro do vídeo. A conferência independente reconstruiu todas as observações, segmentos e janelas a partir dos derivados, verificando 86 arquivos e 9.407.090 comparações. A preparação levou 151,929874 segundos, com pico amostrado de RSS de 226,148 MiB, sem decodificar pixels, avaliar modelos ou abrir fontes de validação e teste. Não houve transição direta de um ID individual para classe 1 nesta referência de treino; a regra foi verificada sinteticamente e não permite inferir a causa biológica de ausências. Em etapa posterior, os índices foram consumidos pelos baselines fixos descritos a seguir, com entradas causais e avaliação densa de ADE.


Como primeira avaliação de predição, persistência e velocidade constante foram fixadas antes dos resultados, sob o commit \texttt{5289c93}. A persistência repete a última posição; a velocidade constante utiliza a mediana componente a componente das cinco últimas diferenças, calculadas a partir de seis posições. Ambos receberam apenas o histórico, em precisão de 64 bits, e produziram todos os passos de 1 a 10 sem limitar as previsões aos contornos da imagem. As 343.776 janelas foram avaliadas por cada método, cobrindo 606 dos 669 IDs individuais originais. Os 63 IDs sem janela permanecem na descrição de cobertura, sem imputação de erro zero. A agregação principal calcula a média de janelas por ID original, reunindo seus segmentos, seguida da média com peso igual entre IDs dentro do vídeo e entre os 12 vídeos. A média por janela dentro de vídeo, seguida de pesos iguais entre vídeos, é secundária.

Na agregação principal, a persistência obteve ADE de 3,890964 pixels e FDE de 6,619751 pixels no horizonte de dez quadros; a velocidade constante obteve 3,137020 e 5,774373 pixels, respectivamente. A redução não foi uniforme: a velocidade constante apresentou ADE menor em nove vídeos e FDE menor em sete. São resultados descritivos nas trajetórias anotadas de treino, condicionados à disponibilidade de futuro de referência, sem ajuste de hiperparâmetros, seleção de método, inferência estatística ou confirmação da hipótese de fluxo. O FPS nominal foi 48 no vídeo 35, 50 no vídeo 82 e 49 nos demais; portanto, dez quadros não correspondem exatamente à mesma duração física em todos os vídeos. Nenhuma reamostragem foi aplicada.

Antes da bateria, 1.362 testes automatizados foram aprovados. A execução completa levou 54,191753 segundos, com RSS amostrado de 199,027 MiB. Uma conferência independente posterior reconstituiu todas as previsões, os erros densos e as agregações a partir dos arquivos exportados e da referência derivada, sem importar a implementação científica ou abrir fontes originais. O manifesto, os arquivos por janela e por ID, a comparação por vídeo e a figura preservam a proveniência da execução. Fluxo, rastreamento e modelos aprendidos não foram avaliados nessa bateria.

Cada trajetória válida será dividida em histórico de 20 quadros e horizontes de 1, 5 e 10 quadros, sem atravessar vídeo, identidade, split ou lacuna de anotação. As entradas básicas serão posição, deslocamento e velocidade estimada. As configurações com fluxo acrescentarão o vetor aparente amostrado na posição da célula ou em um anel válido ao seu redor. Serão comparados persistência, velocidade constante, Filtro de Kalman, Filtro de Partículas e LSTM, além de variantes híbridas correspondentes com fluxo. A LSTM foi escolhida por representar dependências temporais e já ter sido aplicada à previsão da localização de espermatozoides \cite{noy2023location}.

O experimento principal será uma ablação pareada. O mesmo preditor, com os mesmos dados, janelas, hiperparâmetros e sementes, será treinado primeiro sem fluxo e depois com características de fluxo. A variável controlada será, portanto, a presença do movimento aparente, e não uma mudança simultânea de arquitetura. A avaliação ocorrerá em dois cenários: trajetórias anotadas, para isolar o preditor, e trajetórias produzidas pelo rastreador, para medir a degradação da pipeline completa.

Essa ablação exigirá um manifesto comum de janelas, identificadas por vídeo, identidade, instante final do histórico e horizonte. Os dois métodos receberão as mesmas janelas elegíveis; serão informadas a cobertura, as exclusões e as falhas por método. Ausência de fluxo válido não será substituída por vetor zero, nem poderá excluir silenciosamente somente os casos difíceis de um preditor. Executar métodos separadamente não certifica esse pareamento, que será verificado antes de produzir a comparação científica.

Para uma previsão realizada no instante $t$, todas as imagens, máscaras e características de entrada deverão usar somente informações disponíveis até $t$. O fluxo entre $t-1$ e $t$ é admissível; fluxo que utilize imagens posteriores a $t$ não poderá entrar na avaliação principal. O futuro será reservado aos alvos. Qualquer controle que forneça informação futura será identificado como oráculo e relatado separadamente. Normalizações e parâmetros estimados a partir dos dados serão ajustados somente no treino correspondente.


Foi também implementada uma variante causal fixa da velocidade constante,
em precisão de 64 bits. Para as cinco últimas transições observadas, define-se
$\mathbf g_s=F_s(\mathbf p_s)$ e utiliza-se
\begin{equation}
\widehat{\mathbf p}_{t+h}=\mathbf p_t+h\left[
\mathop{\mathrm{mediana}}_{s=t-5}^{t-1}
(\mathbf p_{s+1}-\mathbf p_s-\mathbf g_s)+\mathbf g_{t-1}\right],
\quad h=1,\ldots,10.
\end{equation}
A mediana é calculada por componente. O resíduo é algébrico, sem interpretação
de velocidade intrínseca ou física do fluido. Se o fluxo histórico é constante,
sua subtração e reposição reproduzem o baseline de velocidade constante dentro
da tolerância numérica; esse controle e as recusas de entradas futuras foram
verificados sinteticamente. A avaliação real permanece pendente. Os dois braços
deverão usar a mesma coorte com as cinco amostras finais válidas, reportando
também a cobertura das 19 transições. Ausências anteriores que não entram na
fórmula não excluem por si mesmas a janela, e ausência de fluxo não é imputada
como vetor zero.

\section{Critérios de validação e comparação}

Na detecção, a métrica de promoção será o F1 espacial descrito anteriormente. No rastreamento, a métrica principal será HOTA, acompanhada por MOTA, IDF1, trocas de identidade e fragmentações \cite{luiten2021hota}. Na predição serão calculados ADE e FDE \cite{alahi2016sociallstm}. Os erros serão expressos em pixels e, quando houver calibração confiável, também em micrômetros.

Para cada horizonte $H$, a definição operacional utilizará todos os passos futuros de 1 a $H$:
\begin{equation}
\mathrm{ADE}_{H}=\frac{1}{H}\sum_{h=1}^{H}\left\|\widehat{\mathbf p}_{t+h}-\mathbf p_{t+h}\right\|_2,
\qquad
\mathrm{FDE}_{H}=\left\|\widehat{\mathbf p}_{t+H}-\mathbf p_{t+H}\right\|_2.
\end{equation}
Os horizontes 1, 5 e 10 serão pontos de apresentação; cada ADE exigirá os erros intermediários correspondentes. A média apenas dos erros nos passos 1, 5 e 10 não será denominada ADE de 1 a 10. A interface histórica de avaliação calcula a média dos instantes fornecidos e permanece identificada como tal. A bateria fixa de persistência e velocidade constante utilizou uma interface densa própria, verificada com casos sintéticos e por reconstrução independente dos resultados. As futuras variantes deverão satisfazer o mesmo contrato denso e a separação entre histórico e alvos.

As comparações abrangerão linhas de base clássicas, métodos aprendidos e híbridos, preditores sem e com fluxo e trajetórias anotadas e estimadas. Para evitar que vídeos longos dominem a média, frames e trajetórias serão agregados primeiro dentro de cada vídeo, que será a unidade independente. Friedman para três ou mais métodos e Wilcoxon pareado, com Holm para a família de comparações planejada, somente serão interpretados conforme a adequação do desenho e do tamanho amostral. Serão apresentados os valores e as diferenças pareadas por vídeo, além de média e mediana dos efeitos. Treinamentos estocásticos serão repetidos com sementes distintas, agregadas dentro do vídeo. Qualidade e custo serão apresentados por fronteiras de Pareto, sem um placar único com pesos subjetivos.

No holdout de quatro vídeos, o menor p possível para o Wilcoxon bilateral exato, com diferenças não nulas e sem empates em seus módulos, é $2/2^4=0{,}125$. Assim, esse teste não permite prometer significância a 5\% nesse conjunto. Intervalos de 95\% por bootstrap agrupado de quatro vídeos, caso apresentados, terão caráter exploratório e acompanharão os valores individuais; frames, trajetórias e seeds não ampliarão artificialmente o número de unidades independentes. Uma análise fora da amostra dos 20 vídeos dependerá de seleção e ajuste dentro dos folds de treino de cada avaliação externa, para limitar o viés de seleção\footnote{Cawley e Talbot (2010): \url{https://www.jmlr.org/papers/volume11/cawley10a/cawley10a.pdf}.}. Esse desenho não elimina a necessidade de declarar a exposição histórica do teste.

Os resultados de detecção e rastreamento no VISEM-Tracking serão confrontados com os protocolos publicados por Thambawita et al. \cite{thambawita2023visemtracking} e Zhang et al. \cite{zhang2024spermyolo}. A predição será comparada à linha de base de velocidade constante e ao uso de LSTM descrito por Noy et al. \cite{noy2023location}. Diferenças de divisão dos dados, resolução, horizonte, hardware e implementação serão explicitadas para evitar comparações diretas entre condições incompatíveis.

\section{Etapas de desenvolvimento}

A estratégia inicial priorizou uma pipeline mínima para isolar a hipótese, com referência individual, baselines de persistência e velocidade constante e conexão causal do fluxo. Esses baselines e os derivados de fluxo já foram produzidos e conferidos nas etapas descritas; a ablação real ainda não ocorreu. Na continuidade de 11 de setembro, a comparação e a validação das alternativas de detecção e rastreamento passaram a compor uma frente experimental explícita, iniciada pela busca da Seção~\ref{sec:busca-classicos-20260911}. O refinamento clássico e a preparação autenticada do dataset YOLO foram concluídos, conforme a Seção~\ref{sec:refinamento-classicos-20260911}. As finalistas clássicas ainda exigem validação completa; YOLO precisa de treinamento e comparação no protocolo atual. O contrato de rastreamento deve ser concluído antes da avaliação por HOTA, distinguindo associação com caixas anotadas sem IDs e associação com detecções automáticas.

A escolha da combinação para rastreamento e predição deverá decorrer dessas comparações e da avaliação da cadeia completa, preservando seus erros de localização, identidade e cobertura. Em paralelo, a ablação exploratória nas janelas GT com fluxo já extraído poderá isolar a contribuição dessa característica, sem ser apresentada como resultado do sistema automático. Os modelos aprendidos e seus controles de capacidade deverão ser avaliados antes da síntese final. Um resultado sem redução de ADE/FDE também será informativo, desde que a comparação seja causal, pareada e reproduzível. O roteiro a seguir representa a sequência global do projeto; as etapas já concluídas não serão reiniciadas.

\begin{enumerate}
	\item \textbf{Auditoria e padronização dos dados.} Os vídeos, metadados e anotações serão inventariados, convertidos para formatos comuns e verificados visualmente.

	\item \textbf{Definição das divisões e linhas de base.} Serão fixados os conjuntos de treinamento, validação e teste, as sementes e as configurações clássicas de referência.

	\item \textbf{Treinamento e avaliação da detecção.} A abordagem clássica e o detector aprendido serão ajustados e comparados sob os mesmos dados.

	\item \textbf{Implementação do rastreamento.} Os métodos por centroide, SORT e ByteTrack serão aplicados às mesmas detecções para separar falhas de localização e associação.

	\item \textbf{Estimação e análise do fluxo.} Farnebäck e, se viável, RAFT serão executados com compensação de movimento global e máscaras de primeiro plano.

	\item \textbf{Construção dos conjuntos de predição.} As trajetórias serão segmentadas por continuidade e elegibilidade documentadas, preservando as exclusões. Janelas observadas e futuras serão separadas; normalizações dos preditores serão ajustadas somente no treino correspondente, sem filtrar trajetórias por velocidade ou erro futuro.

	\item \textbf{Treinamento e ablação dos preditores.} As configurações com e sem fluxo serão avaliadas nos cenários controlados.

	\item \textbf{Análise e documentação.} Os resultados serão agregados, acompanhados por intervalos de confiança, visualizações e exemplos de falha, e relacionados à hipótese de pesquisa.
\end{enumerate}

\section{Declaração sobre o uso de Inteligência Artificial}

Para a consolidação desta versão inicial, utilizei as ferramentas ChatGPT e Codex, da OpenAI, como apoio à organização dos conteúdos das entregas anteriores, à pesquisa bibliográfica, à revisão linguística, à conferência da coerência entre os capítulos e à formatação em LaTeX. O tema, os objetivos, as decisões metodológicas e a revisão final do texto são de minha responsabilidade, e a submissão será realizada após a concordância do orientador.

	
	% ---
	% Finaliza a parte no bookmark do PDF, para que se inicie o bookmark na raiz
	% ---
	\bookmarksetup{startatroot} 
	% ---
	
	% ---
	% Conclusão da Entrega 5
	% ---
	\chapter[Conclusão]{Conclusão}

Esta versão inicial consolidou as quatro entregas anteriores na estrutura do modelo de monografia da Faculdade de Computação. O conteúdo da E1 foi reorganizado na Introdução, com separação explícita entre motivação, problema, hipótese, objetivos e contribuições esperadas. Essa mudança eliminou a organização por parágrafos do enunciado e tornou mais clara a relação entre a relevância do tema, a questão investigada e os critérios necessários para verificar o objetivo.

A fundamentação da E2 foi revisada para apresentar os conceitos na ordem em que são utilizados pela pipeline: análise seminal assistida por computador, detecção, rastreamento, fluxo óptico, predição, dados e métricas. Os textos de caráter didático solicitados exclusivamente nas entregas anteriores foram removidos. Os trabalhos analisados na E3 passaram a compor um capítulo próprio, sem os resumos de capítulos do livro e sem informações conjunturais que não contribuem diretamente para a comparação científica. A tabela comparativa foi mantida e atualizada para evidenciar a lacuna relacionada à integração entre rastreamento, movimento aparente e predição individual.

O método proveniente da E4 foi ajustado para manter correspondência direta com a hipótese. Foram detalhadas as divisões dos dados, as linhas de base, a ablação com e sem fluxo, os cenários com trajetórias anotadas e estimadas e os critérios de avaliação. Essa revisão busca reduzir ambiguidades, evitar vazamento de informação e distinguir o movimento aparente medido na imagem de uma velocidade física do fluido. A declaração sobre o uso de Inteligência Artificial foi mantida somente no capítulo de Método, conforme solicitado para esta entrega.

\section{Estado do desenvolvimento e continuidade}

Na revisão de 8 de setembro de 2026, o projeto permanecia alinhado à hipótese e não exigia reinício dos experimentos. A organização das fontes, a definição das divisões e a busca de threshold v3 no treino já haviam avançado além do planejamento inicial. A triagem e o refinamento estavam concluídos e conferidos, com T219/o0/c2 e T218/o0/c2 como finalistas de treino. A auditoria geral registrou 3.248 arquivos conferidos, 17.753 verificações e aprovação em 502 testes automatizados, sem incluir as verificações opcionais de aprendizagem. Esses resultados atestam o alcance das verificações executadas, não a validação científica de todos os módulos nem a hipótese sobre fluxo óptico.

Naquela revisão, o próximo marco previsto era comparar as duas finalistas nos quatro vídeos completos de validação, conforme protocolo prospectivo, totalizando 11.700 avaliações de configuração--quadro. Os controles para quantidade exata de quadros, rejeição de anotações malformadas e integridade das entradas ainda precisavam ser concluídos e registrados. Esse marco é histórico: a bateria e sua conferência foram realizadas posteriormente, como descrito a seguir. Os resultados anteriores permanecem preservados e separados dos contratos novos.

Após essa preparação, a validação completa selecionou T218/o0/c2, com F1 macro de 0,658959, contra 0,657819 de T219/o0/c2. As 11.700 avaliações e a conferência independente foram concluídas; a diferença foi pequena e cada configuração teve F1 maior em dois vídeos. Na conclusão específica daquela bateria ainda não havia congelamento. Posteriormente, T218/o0/c2 foi fixada somente para desenvolvimento; isso não liberou teste, folds ou aplicação, nem demonstrou contribuição do fluxo à predição.

\subsection{Atualização experimental de 11 de setembro de 2026}

A busca clássica entre seis famílias foi concluída e conferida, com 43 configurações nos mesmos 576 quadros dos 12 vídeos de treino, totalizando 24.768 avaliações. Blob obteve o maior F1 macro da grade, 0,791005, enquanto a referência T218 obteve 0,775634. Blob teve F1 maior em sete vídeos e T218 em cinco; a ordem média não é uniforme entre vídeos. A grade limitada, o esforço desigual de busca e o ajuste anterior de T218 impedem interpretar essa triagem como superioridade geral. As 11 finalistas retidas ainda não constituem promoção ou escolha final da pipeline. O protocolo, a revisão operacional e o alcance da conferência estão descritos na Seção~\ref{sec:busca-classicos-20260911}.

A preparação das trajetórias individuais e os baselines de predição sobre GT já estão concluídos. O benchmark compacto de fluxo também foi conferido, mas a hipótese permanece aberta: não há ainda ADE/FDE com fluxo ou HOTA real da comparação entre rastreadores. Após o refinamento local e a preparação do dataset descritos a seguir, a próxima sequência compreende validação completa das alternativas clássicas, treinamento e comparação de YOLO, conclusão do contrato de rastreamento e avaliação dos métodos nas mesmas detecções. Depois de comparar e validar as combinações pertinentes, a predição deverá receber trajetórias efetivamente estimadas, preservando trocas de identidade, perdas e cobertura. A ablação controlada sobre GT pode avançar em paralelo; não substitui a avaliação do sistema automático. Todas as entradas continuarão limitadas ao instante da previsão, e ADE incluirá todos os passos de cada horizonte.

A hipótese ainda não foi testada por esses resultados de detecção. A futura análise deverá declarar a exposição histórica do teste e a limitação de apenas quatro vídeos no holdout, sem prometer significância de 5\% pelo Wilcoxon bilateral exato. Diferenças por vídeo e cobertura das janelas serão essenciais à interpretação. A comparação controlada poderá sustentar, limitar ou não sustentar a contribuição do fluxo; qualquer dessas conclusões dependerá da evidência efetivamente obtida.


\subsection{Avanço após o refinamento local}

O refinamento das cinco famílias clássicas foi concluído e conferido em 45 configurações únicas, com 25.920 avaliações nos mesmos quadros de treino. As dez finalistas retidas ainda não constituem uma escolha final. A paridade dos dez pais, os hashes e a reconstrução das métricas apoiam a consistência dos resultados, mas não substituem a validação em vídeos completos ou a avaliação de rastreamento e predição. O dataset YOLO de 23.316 pares de treino e validação também foi materializado e conferido; não houve treinamento nesta etapa. As regras, o alcance e as limitações desses avanços estão na Seção~\ref{sec:refinamento-classicos-20260911}.

Espera-se que os resultados permitam identificar não apenas a melhor configuração média, mas também as situações em que cada módulo falha. A versão final deverá substituir esta descrição de expectativas por resultados experimentais, discussão das limitações, verificação dos objetivos e síntese das contribuições efetivamente obtidas.
	
	% ----------------------------------------------------------
	% ELEMENTOS PÓS-TEXTUAIS
	% ----------------------------------------------------------
	\postextual
	
	% ----------------------------------------------------------
	% Referências bibliográficas
	% ----------------------------------------------------------
	\bibliography{bib/modelo-referencias}
	
	% ----------------------------------------------------------
	% Glossário - opcional
	% ----------------------------------------------------------
	%
	%\glossary
	
	% Apêndices e anexos não são utilizados nesta versão inicial.
	
	
\end{document}
